\documentclass[10pt,a4paper]{article} 

\usepackage{ctex} % 中文支持
\usepackage[top=2.5cm, bottom=2.5cm, left=2.5cm, right=2.5cm]{geometry} % 页边距
\usepackage{amsmath, amssymb} % 数学公式与符号
\usepackage{graphicx, color, url}

%%%%%%%%%%%%%%%%%%%%%%%%%%%%%%%%%%%%%%%%%%%%%%%%%%%%%%
\usepackage{titling}
\setlength{\droptitle}{-2cm} % 标题上移

\author{王立庆（2025级数学与应用数学1-2班，专业选修课）}
\title{《深度学习基础》教学大纲}
\date{开课学期：2027-2028-1}

%%%%%%%%%%%%%%%%%%%%%%%%%%%%%%%%%%%%%%%%%%%%%%%%%%%%%%

\begin{document}

\maketitle

%%%%%%%%%%%%%%%%%%%%%%%%%%%%%%%%%%%%%%%%%%%%%%%%%%%%%%
\section*{使用教材}
\begin{enumerate}\itemsep0em 
\item 克里斯托弗·M.毕晓普 等 著，邹欣、阮思捷、刘志毅、王树良 译，深度学习：基础与概念，人民邮电出版社，2025年5月第1版。
\item Christopher M. Bishop and Hugh Bishop. Deep Learning: Foundations and Concepts. Springer. 2023. 
\end{enumerate}

%%%%%%%%%%%%%%%%%%%%%%%%%%%%%%%%%%%%%%%%%%%%%%%%%%%%%%
\section*{课程主页}

\url{https://gitee.com/liqing-wang-cscs/deep-learning-foundations}

%%%%%%%%%%%%%%%%%%%%%%%%%%%%%%%%%%%%%%%%%%%%%%%%%%%%%%
\section*{参考网站}
\begin{enumerate}\itemsep0em 
\item  \url{https://www.anthropic.com/}
\item  \url{https://www.aliyun.com/}
\item  \url{https://www.deepseek.com/}
\item  \url{https://arxiv.org/}
\item  \url{https://modelscope.cn/}

\end{enumerate}

%%%%%%%%%%%%%%%%%%%%%%%%%%%%%%%%%%%%%%%%%%%%%%%%%%%%%%
\section*{时间地点}
\begin{enumerate}
\itemsep0em 
\item 上课时间地点：待定。
\item 答疑时间地点：待定。
\end{enumerate}

%%%%%%%%%%%%%%%%%%%%%%%%%%%%%%%%%%%%%%%%%%%%%%%%%%%%%%
\section*{课程成绩}
\begin{enumerate}
\item  平时成绩 100\%, 包括课堂练习、课外作业、课程实验、期末考查。
\begin{enumerate}\itemsep0em 
\item[1.1.] 课堂练习10次，共20分。
\item[1.2.] 课外作业10次，共30分。
\item[1.3.] 课程实验2次，共20分。
\item[1.4.] 期末考查1次，共30分。
\end{enumerate}

%\item  期末成绩 60\%, 计算题和应用题。

\end{enumerate}

%%%%%%%%%%%%%%%%%%%%%%%%%%%%%%%%%%%%%%%%%%%%%%%%%%%%%%
\section*{课程描述}

本课程面向数学与应用数学专业学生，系统讲授深度学习的基本理论与核心算法。内容涵盖概率统计的数学基础、神经网络建模、优化方法、卷积与循环网络等，强调数学基础与算法实现的结合。通过理论学习与编程实践，培养学生运用数学思维分析和解决深度学习有关问题的能力。

%%%%%%%%%%%%%%%%%%%%%%%%%%%%%%%%%%%%%%%%%%%%%%%%%%%%%%
\section*{课程目标}

\begin{enumerate}\itemsep0em 

\item 掌握概率分布、似然函数、和信息论的基本知识。
\item 掌握参数方法（多元正态分布等）和非参数方法（核密度等）。
\item 掌握单层网络（线性回归、分类）的基本知识。
\item 掌握二层网络、混合密度网络的数学原理。
\item 掌握梯度下降等优化方法。
\item 掌握反向传播和自动微分法的数学原理。
\item 掌握正则化的原理和实现技术。
\item 掌握卷积网络的原理和编程实现。

\end{enumerate}

%%%%%%%%%%%%%%%%%%%%%%%%%%%%%%%%%%%%%%%%%%%%%%%%%%%%%%
\section*{授课计划}

\begin{table}[ht!]\centering
\begin{tabular}
{|p{0.8cm}|p{1.4cm}|p{8.5cm}|p{2.6cm}|} \hline 
周	&章	        & 内容     & 案例 \\ \hline\hline  
1	&1.1-1.3    & 深度学习的影响、一个例子、机器学习简史  & 正弦函数拟合 \\  \hline
2	&2.1-2.4    & 概率法则、概率密度、高斯分布、密度变换 & 医学筛查 \\  \hline
3	&2.5-2.6    & 信息论、贝叶斯概率      &  \\  \hline
4	&3.1-3.2	& 离散变量、多元高斯分布      & 老忠实喷泉数据 \\  \hline
5	&3.3-3.5	& 周期变量、指数族分布、非参数化方法      &  \\  \hline
6	&4.1-4.3	& 线性回归、决策理论、偏差-方差权衡      &  \\  \hline
7	&5.1-5.2	& 判别函数、决策理论      &  \\  \hline
8	&5.3-5.4	& 生成分类器、判别分类器      &  \\  \hline
9	&6.1-6.2	& 固定基函数的局限性、多层网络	& 鸢尾花数据 \\  \hline
10	&6.3-6.5	& 深度网络、误差函数、混合密度函数	& 机器人运动学 \\  \hline
11	&7.1-7.2	& 错误平面、梯度下降优化	& \\  \hline
12	&7.3-7.4	& 收敛、Adam算法、归一化	& \\  \hline
13	&8.1-8.2	& 梯度计算、自动微分法	& \\  \hline
14	&9.1-9.6	& 归纳偏置、权重衰减、学习曲线、参数共享、残差连接	& \\  \hline
15	&10.1-10.6	& 卷积滤波器、可视化、目标检测、图像分割、风格迁移	& \\  \hline
\end{tabular}
\end{table}

% 1. 绪论：深度学习发展史与数学基础概览
% 2. 线性代数基础：矩阵运算、特征值分解与奇异值分解在神经网络中的应用
% 3. 多元微积分：梯度、雅可比矩阵与海森矩阵的推导及计算图原理
% 4. 概率论与信息论：最大似然估计、贝叶斯推断与交叉熵损失函数
% 5. 优化理论基础：梯度下降法、动量法与自适应学习率算法（Adam）的收敛性分析
% 6. 前馈神经网络：万能近似定理、反向传播算法的数学推导与正则化技术
% 7. 卷积神经网络：卷积运算的数学定义、池化操作与经典架构（ResNet）分析
% 8. 循环神经网络：序列建模、梯度消失与爆炸问题的数学机理及LSTM结构
% 9. 注意力机制与Transformer：自注意力矩阵运算、位置编码与并行化计算优势
% 10. 生成模型基础：变分自编码器的证据下界推导与扩散模型的随机微分方程视角
% 11. 深度强化学习：马尔可夫决策过程、贝尔曼方程与策略梯度定理
% 12. 深度学习的泛化理论：VC维、Rademacher复杂度与神经正切核简介
% 13. 深度学习的逼近理论：深度网络对高维函数的表达能力与效率分析
% 14. 科学计算中的深度学习：利用神经网络求解偏微分方程（PINNs）
% 15. 课程总结与前沿研讨：模型可解释性、鲁棒性与大模型数学原理展望

%%%%%%%%%%%%%%%%%%%%%%%%%%%%%%%%%%%%%%%%%%%%%%%%%%%%%%

\end{document}


